\section{Cluster Dynamics Formulation}\label{sec:formulation}

\subsection{From the master equation to the reduced system}
\label{sec:master}

The general cluster-dynamics master equation for a population indexed by a
species label $\alpha$ and a size $n$ is a balance of gains and losses over the
edges of the graph of \cref{fig:rag},
\begin{equation}
  \frac{\partial C^\alpha_n}{\partial t}
  = -\nabla\!\cdot\!\bm J^\alpha_n
    + G^\alpha_n
    + \sum_{\beta,n'} \bigl[\mathcal K^{\alpha\beta}_{n-n',n'}C^\beta_{n-n'}C^{\beta'}_{n'}
      - \mathcal K^{\alpha\beta}_{n,n'}C^\alpha_nC^\beta_{n'}\bigr]
    + \varepsilon^\alpha_{n+1}C^\alpha_{n+1} - \varepsilon^\alpha_nC^\alpha_n
    - \mathcal D^\alpha_n C^\alpha_n ,
  \label{eq:master}
\end{equation}
with $\bm J^\alpha_n=-\bm D^\alpha_n\nabla C^\alpha_n$ the diffusive flux,
$G^\alpha_n$ the cascade source, $\mathcal K$ the pairwise reaction kernels,
$\varepsilon$ the thermal emission rates, and $\mathcal D$ the absorption rate
at fixed sinks. \Cref{eq:master} is one equation per $(\alpha,n)$ and per mesh
node, and is not solvable as written for the size ranges of interest.

Two reductions bring it to the system integrated here, and they cut in different
directions.

\paragraph{Mobile species: truncate the ladder.} Only the monomers and the
smallest interstitial clusters migrate. The mobile set is closed at
$\spvec{m}_M=[-1,+1,+2,+3]^{T}$ --- monovacancy, mono-, di- and tri-interstitial
--- with the signed entries $m^k$ carrying the polarity
$p^k=m^k/|m^k|$. Everything larger is immobile and is carried by the second
reduction. The mobile block retains its spatial derivative, because transport is
what it is for.

\paragraph{Immobile families: truncate the moments, keep the variants.} For the
immobile families the ladder in $n$ is replaced by the first three moments of
the size distribution of each family in \cref{eq:families},
\begin{equation}
  n^k_{sL}=\!\int\! f_k\,\mathrm{d}m,
  \qquad
  c^k_{sL}=\!\int\! m f_k\,\mathrm{d}m,
  \qquad
  q^k_{sL}=\!\int\! m^2 f_k\,\mathrm{d}m,
  \label{eq:moments-def}
\end{equation}
closed by the log-normal ansatz of \cref{sec:moments}. The immobile block loses
its spatial derivative entirely: a loop does not diffuse, and the only thing
that couples one node's loops to another's is the mobile field. That is not an
approximation --- it is the property that makes the numerical method of
\cref{sec:split} possible.

\subsection{The mobile block}
\label{sec:mobile}

Because the mobile species reach quasi-steady state on a time scale
$\tau_M\sim(k^2D)^{-1}$ of order microseconds, while the immobile families evolve
on the dose scale ($1\,\mathrm{dpa}=10^{7}\,\mathrm{s}$ at
$G=10^{-7}\,\mathrm{dpa\,s^{-1}}$), the mobile block is solved as a
\emph{steady} reaction--diffusion problem for the immobile state currently held.
Writing $\spvec c_M=[c^v,c^i,c^{2i},c^{3i}]^{T}$ in per-atom units,
\begin{equation}
  \spvec 0 = -\nabla\!\cdot\!\spvec J
    + \spvec P_M
    + \tfrac12\,(\spmat I\otimes\spvec c_M^{T})\,\spmat R_2\,\spvec c_M
    + (\spmat E-\spmat D)\,\spvec c_M
    + \spvec S_{vL},
  \label{eq:master-c}
\end{equation}
where upright sans-serif bold denotes arrays in \emph{species space} and
bold italic ordinary Cartesian vectors and tensors, so that the entries
$\bm J^{x}=-\bm D^{x}\nabla c^{x}$ of the stacked flux $\spvec J$ are themselves
Cartesian vectors and $\nabla\cdot$ acts on each.

The terms are, in order: the cascade production
$\spvec P_M=G\varepsilon\spvec f$, with $G$ the displacement rate, $\varepsilon$
the surviving efficiency and $\spvec f$ the partition of production among the
mobile species; the second-order reactions $\spmat R_2$, whose entries are the
pairwise rate constants
\begin{equation}
  \beta^{m,n}=\frac{A^{m,n}}{\Omega}\,(\overline D^{m}+\overline D^{n})(r^{m}+r^{n}),
  \qquad
  r^{m}=\begin{cases}
    \left(\dfrac{3\Omega}{4\pi}\right)^{1/3}, & |m^{m}|=1,\\[9pt]
    \left(\dfrac{|m^{m}|\,\Omega}{\pi b}\right)^{1/2}, & |m^{m}|>1,
  \end{cases}
  \label{eq:rates}
\end{equation}
with the recombination-volume factor $\beta_r$ carried on the
interstitial--vacancy channel and the combinatorial factors retained on the
like-species channels; the thermal emission matrix $\spmat E$, whose entries are
$\alpha^{m,n}=\beta^{m,n}\exp(-E^{b}_{m,n}/k_BT)$; the linear absorption matrix
$\spmat D$; and $\spvec S_{vL}$, the vacancies returned to the pool by
peripheral emission from vacancy loops (\cref{app:emission}).

The absorption matrix is where the immobile population acts back on the mobile
one, and it is written so that no defect can be counted twice:
\begin{equation}
  \spmat D = \mathrm{diag}\{\spvec\omega\odot\spvec C_s\}
             + \spmat D_d\,\spmat K_{\mathcal L},
  \qquad
  \spmat D_d=\mathrm{diag}\{\overline D^{v},\overline D^{i},
                            \overline D^{2i},\overline D^{3i}\},
  \label{eq:KM}
\end{equation}
whose first term carries the \emph{network alone},
\begin{equation}
  C_s^{v}=\frac{a^2}{z_c}\rho_N,
  \qquad
  C_s^{i}=\frac{a^2}{z_c}Z_N\rho_N,
  \label{eq:sinkconc}
\end{equation}
and whose second carries the loops through the diagonal sink-strength block
\begin{equation}
  \spmat K_{\mathcal L}=\mathrm{diag}\{\spvec k^{2}\},
  \qquad
  k_m^{2}=\sum_{(s,k)\in\mathcal F_L}\! Z_{sk,m}\,\rho^k_{sL}
        \;+\;Z_{vc_0,m}\,\alpha_{\mathrm{sfp}}\,\pi R^{c_0}_{vL}\,
             \frac{n^{c_0}_{vL}}{\Omega}
        \;+\;Z_{N,m}\,\rho_N .
  \label{eq:KL}
\end{equation}
The loops must not also appear in \cref{eq:sinkconc}, or their absorption would
be counted once as a network sink and once as loop growth. The middle term of
\cref{eq:KL} is the stacking-fault pyramid, which has no line density and enters
through a compact capture cross-section instead; \cref{eq:Pi-identity} below
shows that a cross-section is the only thing either side ever asks of a family.

\paragraph{Boundary conditions.} On a free surface or a grain boundary the
mobile species are pinned at thermal equilibrium,
\begin{equation}
  c^{v}=c^{v,\mathrm{eq}}=\exp(-E^{f}_v/k_BT),
  \qquad
  c^{i}=c^{2i}=c^{3i}=0 \quad\text{on }\partial\Omega_D,
  \label{eq:bc}
\end{equation}
which is process class P13 of \cref{tab:rag-processes}. This replaces the
mean-field grain-boundary sink strength $k^2_{gb}\simeq6\sqrt{k^2}/d$ by a
boundary condition, and the depleted layer in front of it becomes an output
rather than a parameter. Faces declared periodic instead carry periodic
conditions on all four species.

\subsection{Anisotropic transport and the single anisotropy parameter}
\label{sec:transport}

Each mobile species carries a diffusion tensor
$D^{x}_{ij}=D^{x}_{0ij}\exp(-E^{xm}_{ij}/k_BT)$, diagonal in the crystal frame
with the basal components equal and the $c$-axis component distinct. The whole
of that anisotropy is expressed by one number per species,
\begin{equation}
  p^{m}=\left(\frac{D^{m}_{\parallel}}{D^{m}_{\mathrm{basal}}}\right)^{1/6},
  \label{eq:pm}
\end{equation}
and the migration energies are generated from it at \emph{fixed} effective
diffusivity,
\begin{equation}
  E^{m}_{\langle11\rangle}=E^{m}_{\langle22\rangle}=E^{m}_{\mathrm{eff}}+2k_BT\ln p^{m},
  \qquad
  E^{m}_{\langle33\rangle}=E^{m}_{\mathrm{eff}}-4k_BT\ln p^{m},
  \label{eq:Emsplit}
\end{equation}
so that the orientation-averaged diffusivity
$\overline D^{m}=(\det\bm D^{m})^{1/3}$ is exactly the scalar $D_m$ of the
mean-field kinetics. \textbf{No mobility is redefined by the anisotropy --- only
its directional weighting.} This is the technical device that makes the
formulation self-consistent: \cref{eq:pm,eq:Emsplit} are a change of variable in
which $p^{m}$ is the only anisotropy the model has, and \cref{sec:efficiencies}
uses that same $p^{m}$ and no other parameter.

\subsection{Capture efficiencies}
\label{sec:efficiencies}

Every capture efficiency is a product of factors, one per degeneracy to be
broken:
\begin{equation}
  Z_{sk,m}(\bm\sigma)
  = \underbrace{Z^{0}_m}_{\text{elastic bias}}\;
    \underbrace{A_{h(k)}(p^{m})}_{\text{DAD, orientation}}\;
    \underbrace{X_{s,m}}_{\text{character}}\;
    \underbrace{\left[1+\frac{\Delta Z^{\mathrm{SIPA}}_{h(k),m}(\sigma_{kk})}
                            {Z^{0}_mA_{h(k)}(p^{m})}\right]}_{\text{stress, SIPA}},
  \label{eq:Zsk}
\end{equation}
with $h(k)\in\{a,c\}$ the habit of family $k$ and the Woo diffusional-anisotropy
factors~\citep{woo1988theory,woo1981sink,po2026tensorialdad}
\begin{equation}
  A_a(p)=\frac{p+p^{-2}}{2},
  \qquad
  A_c(p)=p .
  \label{eq:Afactors}
\end{equation}
Each factor breaks exactly one degeneracy and is silent on the others.
$A_{h(k)}$ depends on the habit alone --- on the angle between the loop normal
and the $c$ axis --- and therefore assigns one value to all six prismatic
families and one to all three basal states, the pyramid included, whose base is
$(0001)$ so that $h(c_0)=c$. It separates prismatic from basal but is degenerate
in character. $X_{s,m}$ is the character factor, which separates an interstitial
from a vacancy loop of the same habit~\citep{marchrico2023,marchrico2025}; the
coexistence window it opens has log-width $\ln\chi$ with
$\chi=X_{iI}X_{vV}/(X_{iV}X_{vI})$.

\textbf{\Cref{eq:Zsk} contains no free bias.} $A_{h(k)}$ is built from the same
$p^{m}$ that \cref{eq:Emsplit} puts into the diffusion tensors, so a change in
the transport anisotropy moves the growth laws and the sink strengths of
\cref{eq:KL} together. The two remaining constants, $Z^{0}_m$ and $X_{s,m}$, are
calibrated (\cref{sec:calibration}); the orientation dependence is not.

The forms in \cref{eq:Afactors} are not monotone in the same way, and the
consequence is quantitative and easily missed. $A_c(p)=p$ increases without
bound, while $A_a(p)=(p+p^{-2})/2$ has its minimum at $p=2^{1/3}=1.2599$. A
single change of $p^{v}$ therefore pushes the basal and prismatic families in
\emph{opposite} directions, starving \cloop{} growth while accelerating \aloop{}
shrinkage, or the reverse. Any statement about the effect of anisotropy that
does not distinguish the two habits is incomplete.

\subsection{Conditions for the growth of every variant}
\label{sec:coexistence}

The experimental record poses a problem that a mean-field model does not
automatically solve, and the analysis that follows is what forces the
formulation to be self-consistent in the sense of
\cref{sec:intro}(a). It is also the analysis that decides which
of the nine families grow, at what mobile fluxes, and under what conditions two
of them can grow at once.

\subsubsection{The growth conditions}
\label{sec:growconds}

Define the mobile flux aggregates and their ratio,
\begin{equation}
  \Phi_I=\sum_{n\in I_{\mathrm{mob}}}n\,\overline D^{I_n}c^{I_n},
  \qquad
  \Phi_V=\overline D^{v}c^{v},
  \qquad
  x\equiv\frac{\Phi_V}{\Phi_I},
  \label{eq:fluxscales}
\end{equation}
which reduce to $\Phi_I=\overline D^{i}c^{i}$ and
$\Phi_V=\overline D^{v}c^{v}$ in the monomer approximation. With $Z_{s,m}$ the
capture efficiency of family $s$ for mobile species $m$, the net driving force
on a family of character sign $\eta_s=\pm1$ is
\begin{equation}
  F_s=\eta_s\left(Z_{s,I}\Phi_I-Z_{s,V}\Phi_V\right),
  \label{eq:Fsigma}
\end{equation}
so that each family grows or shrinks according to which side of its own
threshold the ratio $x$ falls on:
\begin{equation}
  \text{interstitial family }s\text{ grows}\iff x<\frac{Z_{s,I}}{Z_{s,V}},
  \qquad
  \text{vacancy family }s\text{ grows}\iff x>\frac{Z_{s,I}}{Z_{s,V}} .
  \label{eq:growconds}
\end{equation}

\Cref{eq:growconds} is the whole of the growth problem, and it says something
uncomfortable. Every family reduces to \emph{one} threshold, the ratio of its
own two efficiencies, and $x$ is a single number. Under identical local fluxes
and identical capture coefficients an interstitial and a vacancy loop of the
same orientation therefore \emph{cannot} both grow: whichever side of the shared
threshold $x$ lies on, one of them is shrinking. Yet both prismatic characters
are observed together in irradiated
Zr~\citep{jostsons1977nature,griffiths1988review,carpenter1981study,%
marchrico2025}. Coexistence is thus a statement that at least one idealization
has failed --- homogeneity, identical capture coefficients, identical nucleation
pathways, identical variants, or identical irradiation history --- and the
remainder of this subsection identifies which, and what each candidate can and
cannot deliver.

\subsubsection{What diffusional anisotropy does, and where it stops}
\label{sec:dad-window}

With the diffusion tensors and the orientation factors of
\cref{eq:pm,eq:Afactors}, the prismatic and basal thresholds separate:
\begin{equation}
  x_a=\frac{Z^{0}_{I}A_a(p^{i})}{Z^{0}_{V}A_a(p^{v})},
  \qquad
  x_c=\frac{Z^{0}_{I}A_c(p^{i})}{Z^{0}_{V}A_c(p^{v})},
  \label{eq:thresholds}
\end{equation}
and because $A_a(p)=(p+p^{-2})/2$ exceeds $A_c(p)=p$ for $p<1$ while the two
coincide at $p=1$, the ordering $x_c<x_a$ holds whenever $p^{i}<p^{v}$.
Prismatic interstitial and basal vacancy loops then grow \emph{together}
throughout
\begin{equation}
  \boxed{\;x_c<x<x_a\;},
  \qquad
  \ln\frac{x_a}{x_c}=\ln\frac{g(p^{v})}{g(p^{i})},
  \qquad
  g(p)\equiv\frac{2}{1+p^{-3}},
  \label{eq:window}
\end{equation}
with $g$ strictly increasing, so the window is open exactly when $p^{i}<p^{v}$
and closes as $p^{i}\to p^{v}$. In the near-isotropic-vacancy limit
$p^{v}\simeq1$ the width reduces to the familiar form
$x_a-x_c=(Z^{0}_{I}/Z^{0}_{V})(1-(p^{i})^{3})/2(p^{i})^{2}$. At
$200\,^{\circ}$C, with $Z^{0}_{I}/Z^{0}_{V}\approx1.05$--$1.15$ and
$p^{i}\approx0.89$--$0.92$~\citep{woo1988theory,christien2009cluster,%
barashev2015theoretical,li2020coupled}, the window is open but narrow:
$x_c\approx1.0<x<x_a\approx1.2$.

Three things are worth separating here, because they are routinely conflated.

\textbf{The elastic bias opens nothing.} $Z^{0}_{I}/Z^{0}_{V}>1$ multiplies
\emph{both} thresholds in \cref{eq:thresholds} and so shifts them together; it
cancels from the width in \cref{eq:window} entirely. It does narrow the
prismatic-interstitial/basal-vacancy window in a practical sense, by pushing
$x_c$ toward unity, but it cannot create an interval where none exists.

\textbf{What positions $x$ inside the window is the production bias.} The
excess of freely migrating vacancies left when more of the interstitial yield is
locked into in-cascade clusters, $\eta=\varepsilon_V/\varepsilon_I\gtrsim1$,
sets where in $(x_c,x_a)$ the operating point sits; a $0$--$20\%$ excess
suffices. \textbf{Neither substitutes for the other.} DAD opens the interval and
the production bias places the point inside it: at $p^{i}=p^{v}$ the thresholds
coincide and no value of $\eta$ can satisfy both conditions at once, while with
the window open but $\eta$ badly placed the model simply grows one family and
dissolves the other.

\textbf{DAD is degenerate in character, and that is where it stops.}
\Cref{eq:Zsk} assigns one orientation factor $A_{h(k)}(p^{m})$ to every family
sharing a habit, so the interstitial and the vacancy prismatic variants receive
the \emph{same} pair of efficiencies and hence the same threshold,
\begin{equation}
  \frac{Z_{a_k I,I}}{Z_{a_k I,V}}
  \;\xrightarrow[\ \text{DAD alone}\ ]{}\;
  \frac{Z_{a_k V,I}}{Z_{a_k V,V}} .
  \label{eq:dad-degenerate}
\end{equation}
The prismatic window closes identically, and DAD is silent on the coexistence of
interstitial and vacancy \aloop{} loops on one habit --- which is the
coexistence the experiments report. Two mechanisms remove that degeneracy, and
the formulation carries both.

\subsubsection{Resolution 1: the flux ratio is a field}
\label{sec:resolution-field}

Promote $x$ to $x(\xv,t)=\Phi_V(\xv,t)/\Phi_I(\xv,t)$. Because interstitials
migrate far faster than vacancies, the interstitial channel responds first and
most sharply to any transient or gradient.

In \emph{time}, the nucleation transient --- which in Zr extends over of order
$1\,$dpa --- is strongly interstitial-rich, so $x(t)\ll x_c$ and interstitial
loops grow while vacancy loops merely nucleate and sit sub-critical; as the
vacancy supersaturation builds, $x(t)$ rises into $(x_c,x_a)$ and the vacancy
families begin to grow, with the interstitial loops established earlier
persisting alongside them.

In \emph{space}, sinks --- grain boundaries, free surfaces, precipitates, the
network --- make the mobile fields non-uniform, and because the interstitial
diffusion length far exceeds the vacancy one, interstitials drain into the far
field while vacancies pile up locally, so $x(\xv)$ is depressed far from sinks
and enhanced near them. A single specimen then contains regions on both sides of
the prismatic threshold. \Cref{sec:results-mobile} measures exactly this
structure: boundary layers of $8$--$46\,$nm that differ by species, so $x$
varies systematically across them.

\textbf{This mechanism costs nothing to include: a spatially resolved
formulation carries $x(\xv,t)$ natively and does not have to be told to.} Its
limitation is equally plain. It leaves $Z_{a_kI,\cdot}=Z_{a_kV,\cdot}$ intact
and buys coexistence by declining to evaluate the two growth conditions at the
same point of space and time, so it predicts \emph{segregated} coexistence and
cannot give locally simultaneous growth in a homogeneous volume.

\subsubsection{Resolution 2: capture is character-resolved}
\label{sec:resolution-character}

March-Rico, Blondel and Wirth~\citep{marchrico2025} compute the defect--loop
interaction energy atomistically and replace geometric intersection by a
drift--diffusion capture,
\begin{equation}
  \bm J_\alpha=-\bm D_\alpha\left[\nabla c^{\alpha}
    +\frac{c^{\alpha}}{k_BT}\nabla E^{\alpha L}_{\mathrm{int}}\right],
  \qquad \alpha=i,v,
  \label{eq:drift-diffusion}
\end{equation}
which reduces in the rate equations to effective capture radii entering
additively in the reaction volume. \emph{Spontaneous} capture radii are nearly
the same for all loops and enlarge every reaction volume without separating
characters; \emph{thermal-drift} radii do separate them, because the dilatational
field of an interstitial loop attracts interstitials over a longer range than it
attracts vacancies, and conversely for a vacancy loop. The result is a
\textbf{same-type capture bias}, carried in \cref{eq:Zsk} by the factor
$X_{s,m}$. Defining the character-splitting factor
\begin{equation}
  \chi\equiv
  \frac{Z_{a_kI,I}/Z_{a_kI,V}}{Z_{a_kV,I}/Z_{a_kV,V}}
  =\frac{X_{iI}X_{vV}}{X_{iV}X_{vI}},
  \label{eq:chi}
\end{equation}
the prismatic window $Z_{a_kV,I}/Z_{a_kV,V}<x<Z_{a_kI,I}/Z_{a_kI,V}$ is
non-empty if and only if $\chi>1$, with logarithmic width $\ln\chi$.

\Cref{eq:chi} is the exact character analogue of the orientation width
\cref{eq:window}: $\ln(x_a/x_c)$ measures how far diffusional anisotropy
separates prismatic from basal, and $\ln\chi$ how far thermal-drift capture
separates interstitial from vacancy \emph{within} one prismatic habit. Because
$\chi>1$ opens a genuine interval, a single uniform steady $x$ satisfies both
conditions at once --- which is the stronger claim, and the more falsifiable
one.

\subsubsection{Two ingredients that are not optional}
\label{sec:not-optional}

\textbf{In-cascade vacancy clustering supplies the embryos.} Without them the
capture bias has nothing to act on: homogeneous nucleation of a vacancy loop
against the prevailing interstitial bias is prohibitively slow. It is the
\emph{same} clustering asymmetry that generates the production bias $\eta$, so
the two must not be fitted as independent inputs.

\textbf{Cascade overlap reduces the effective generation rate with dose,}
$\tilde g(\gamma)=\varepsilon(\gamma)g$, saturating over $0.01$--$0.1\,$dpa to a
small residual fraction. This is what lets coarsening rather than unbounded
nucleation dominate at high dose. Applied as a single scalar to both species it
rescales dose without moving $x$: it is a rate effect, not a threshold effect.

\Cref{tab:mechanisms} summarizes the division of labor.

\begin{table}[htbp]
\caption{What each ingredient does to the growth criterion \cref{eq:growconds}.
Only the last two separate loops of the same prismatic orientation; only
character-resolved capture does so locally and instantaneously.}
\label{tab:mechanisms}
\centering\small
\begin{tabularx}{\textwidth}{@{}>{\raggedright\arraybackslash}p{0.20\textwidth}
                              >{\raggedright\arraybackslash}p{0.20\textwidth}X@{}}
\toprule
Ingredient & Acts on & What it does, and does not do \\
\midrule
Elastic bias $Z^{0}_{I}/Z^{0}_{V}$ & both thresholds together &
  Shifts $x_a$ and $x_c$ jointly; opens no window, and narrows the
  prismatic-interstitial/basal-vacancy window by pushing $x_c$ toward unity. \\
\addlinespace
Diffusional anisotropy, $p^{i}<p^{v}$ & orientation &
  Opens the prismatic-interstitial/basal-vacancy window of width
  \cref{eq:window}. Cannot separate loops of one orientation,
  \cref{eq:dad-degenerate}. \\
\addlinespace
Production bias $\eta$ & position of $x$ &
  Places $x$ inside whatever window exists. Creates no window; with $x_a=x_c$ it
  merely switches which family grows. \\
\addlinespace
Flux field $x(\xv,t)$ & position of $x$, resolved in space and time &
  Lets the two prismatic conditions be met at different doses or different
  places. Predicts \emph{segregated} coexistence. Native to a spatially
  resolved model. \\
\addlinespace
Thermal-drift capture, $\chi>1$ & character &
  Opens the prismatic window of logarithmic width $\ln\chi$; gives \emph{local}
  simultaneous growth. Needs in-cascade clustering for the embryos, and still
  needs $x$ positioned inside. \\
\bottomrule
\end{tabularx}
\end{table}

\subsubsection{Why this forces the formulation}
\label{sec:forces-formulation}

Three requirements follow, and together they are the content of
\cref{sec:classes,sec:formulation}.

\begin{enumerate}[itemsep=3pt]
\item \textbf{The immobile families must be crystallographic.} A scalar
      aligned/non-aligned split can express \emph{how much} of a population a
      stress favors but not \emph{which} prismatic system it favors, and it has
      no meaning at all under a multiaxial stress or a stress field. Resolving
      $a_1,a_2,a_3$ makes the resolved stress on each habit a computed
      quantity, and carrying both characters on each variant is what makes
      \cref{eq:chi} representable at all. Splitting the basal population into a
      faulted and a perfect state is forced separately, by the fact that the two
      store different vacancy content per unit area (\cref{tab:cd-families}).
\item \textbf{The capture efficiencies must come from the transport tensors.}
      This is sense~(a) of \cref{sec:intro}, and
      \cref{eq:window} is why it matters rather than merely being tidy: the
      width of the co-growth window is a function of the transport anisotropy
      \emph{alone}. A model that diffuses with one anisotropy and biases its
      sinks with an independent phenomenological pair
      $(\delta_i,\delta_v)$ can therefore report a window that its own transport
      does not support. The phenomenological pair remains recoverable as the
      exact map
      \begin{equation}
        \delta_m=\frac{A_a(p^{m})-A_c(p^{m})}{A_a(p^{m})+A_c(p^{m})}
                =\frac{1-(p^{m})^{3}}{1+3(p^{m})^{3}},
        \label{eq:deltamap}
      \end{equation}
      which gives $\delta_i\approx0.07$--$0.09$ and $\delta_v\approx0$ at the
      $200\,^{\circ}$C anisotropies above. Note that $Z^{0}_{I}/Z^{0}_{V}$ does
      not appear in $\delta_m$ and must be carried separately, which is why
      \cref{eq:Zsk} is the more faithful input.
\item \textbf{The mobile field must be spatially resolved.}
      \Cref{sec:resolution-field} is not available to a $0$-D model. It requires
      $x(\xv,t)$, hence the reaction--diffusion problem of \cref{eq:master-c}
      with real boundary conditions --- and it is one of the two independent
      reasons for the spatial treatment, the other being that surfaces and grain
      boundaries are the dominant defect sink in any specimen small enough to be
      examined (\cref{tab:conservation}).
\end{enumerate}

\paragraph{Two cautions, so that neither is later mistaken for a prediction.}
The thermal-drift radii behind \cref{eq:chi} scale as $(T_m/T)^{m_d}$ with
$T_m=2128\,$K, so with $m_d>0$ the same-type bias, and hence $\chi$, is
\emph{strongest at low temperature} --- while the observed vacancy \aloop{}
fraction \emph{increases} with temperature. These are not contradictory: the
temperature dependence of the vacancy fraction is controlled largely by
nucleation, by vacancy mobility, and by thermal emission from small vacancy
clusters, none of which is contained in $\chi$. But reproducing the measured
vacancy fraction against temperature is a genuine test of the combined model
rather than something it obtains for free. The related hazard is double
counting: the same-type capture bias \emph{is} an elastic-interaction effect,
and so is the conventional dislocation bias $Z^{0}_{I}/Z^{0}_{V}>1$. Only one of
the two should carry the interstitial--vacancy asymmetry. In the calibration of
\cref{sec:calibration} the fitted $\chi=1$, so the asymmetry is carried by
$Z^{0}$ alone and the question does not arise --- which is also to say that the
present parameter set exercises Resolution~1 and not Resolution~2.

\subsection{The immobile block}
\label{sec:immobile}

\paragraph{Growth.} All nine families obey one growth law, in which the stored
content changes under the net biased flux of the correct sign:
\begin{equation}
  \left.\dot c^k_{sL}\right|_{\mathrm{growth}}
  = \Pi^k_{sL}\;\eta_s\!\!\sum_{m\in\{v,i,2i,3i\}}\!\!
    |m|\;\overline D^{m}\,Z_{sk,m}\,c^{m}\,\varsigma_s(m),
  \label{eq:growth}
\end{equation}
with $\varsigma_s(m)=+1$ when the polarity of species $m$ matches the character
$s$ and $-1$ otherwise. The geometric prefactor is
\begin{equation}
  \Pi^k_{sL}=\frac{\lambda_k}{\ell}\,\varsigma^{\mathrm{sink}}_k
             \sqrt{n^k_{sL}c^k_{sL}},
  \qquad
  \ell\equiv\frac{z_c\Omega}{2\pi a^{2}},
  \label{eq:loopprefactor}
\end{equation}
and because $\rho^k_{sL}=2\pi\lambda_k\sqrt{n^k_{sL}c^k_{sL}}$ by
\cref{eq:radius},
\begin{equation}
  \Pi^k_{sL}=\varsigma^{\mathrm{sink}}_k\,\frac{\rho^k_{sL}}{2\pi\ell}
  \qquad\text{identically, for every }(s,k)\in\mathcal F_L .
  \label{eq:Pi-identity}
\end{equation}
\emph{The prefactor in the growth law is the family's own sink density divided
by $2\pi\ell$.} That identity is what makes \cref{eq:growth} and \cref{eq:KL}
two readings of one quantity, and it is what carries the growth law onto the
pyramid unchanged: a family needs no geometry beyond the capture cross-section
it presents, and the pyramid supplies that cross-section as
$\Pi^{c_0}_{vL}=\varsigma^{\mathrm{sink}}_{c_0}\alpha_{\mathrm{sfp}}
R^{c_0}_{vL}n^{c_0}_{vL}/2\ell$.

\paragraph{Nucleation.} Cascade production deposits directly into the immobile
families at a rate $G_k=G\,\epsilon_k$ per family, each new cluster carrying
$m^{\mathrm{nuc}}_k$ defects, so that
$\dot n^k_{sL}|_{\mathrm{nuc}}=J^k_{sL}=G_k/m^{\mathrm{nuc}}_k$ and
$\dot c^k_{sL}|_{\mathrm{nuc}}=G_k$. The partition $\epsilon_k$ and the
nucleation sizes are calibrated quantities.

\paragraph{Coalescence.} Loops are removed by overlap with other loops and with
the network. With $N=n/\Omega$, $R=\lambda\sqrt{c/n}$, and the radii capped at
the relevant spacing ($R_{LL}=\min(R,(\Omega/n)^{1/3})$,
$R_{LN}=\min(R,\rho_N^{-1/2})$), the Avrami overlap probabilities and rate
scales are
\begin{align}
  \varphi_{LL}&=1-\exp\!\left(-\kappa_{LL}\tfrac43\pi R_{LL}^{3}N\right), &
  \varphi_{LN}&=1-\exp\!\left(-\kappa_{LN}\pi R_{LN}^{2}\rho_N\right),
  \nonumber\\
  \nu_{LL}&=c_{LL}\,v_{\mathrm{abs}}\,N^{1/3}, &
  \nu_{LN}&=c_{LN}\,v_{\mathrm{abs}}\,\sqrt{\rho_N},
  \label{eq:avrami}
\end{align}
giving the number and content losses
\begin{equation}
  \Delta N_{LN}=\nu_{LN}\varphi_{LN}n,
  \qquad
  \Delta N=\nu_{LL}\varphi_{LL}n+\Delta N_{LN},
  \qquad
  \Delta C=\nu_{LN}\varphi_{LN}c .
  \label{eq:coal}
\end{equation}
The like-loop channel merges two loops into a larger one and therefore
\emph{conserves content}, removing density only; the loop--network channel
removes both, and the removed content is booked as network sink absorption so
the point-defect balance stays closed. The rate scale is set by the
\emph{absorbed}-flux climb speed $v_{\mathrm{abs}}$ rather than the net growth
velocity, so coarsening persists at saturation. The capping is what makes
\cref{eq:avrami} saturate rather than diverge as the number density collapses,
and $\varphi_{LL}\to1$ is the model's own statement that its mean-field
treatment of coalescence has expired --- the detector used in
\cref{sec:coarsening}.

\paragraph{Emission.} Vacancy loops exchange vacancies across their own
perimeter at a rate set by a size-dependent binding energy computed from the
anisotropic loop self-energy (\cref{app:emission}), one expression serving both
characters, so that the same mechanism shrinks a vacancy loop and grows an
interstitial one and the classical climb threshold emerges rather than being
imposed. This channel replaces the two fitted annealing lifetimes of the
mean-field model. It returns vacancies to the mobile pool through $\spvec S_{vL}$
in \cref{eq:master-c}.

\paragraph{The basal chain.} The three basal states are linked by two one-way
activated transfers,
\begin{equation}
  c_0 \;\xrightarrow{\;\nu_{\mathrm{col}}\;}\; c_f
      \;\xrightarrow{\;\nu_{\mathrm{uf}}\;}\; c_p ,
  \label{eq:basal-chain}
\end{equation}
both size-gated, and both conserving vacancies exactly: the content leaving one
state enters the next, and only the \emph{number} changes, because the states
store different vacancy content per unit area
(\cref{tab:cd-families})~\citep{liu2020two,varvenne2014vacancy,marchrico2025}.
The chain is basal and has no prismatic counterpart: \emph{ab initio}-parameterized
formation energies place the prismatic vacancy \aloop{} loop \emph{below} both
the compact cavity and the faulted basal loop, so the planar prismatic loop is
already the ground-state geometry rather than something a compact cluster must
convert into~\citep{varvenne2014vacancy}. \textbf{None of the basal-chain rates
is calibrated}; the formulation gives their Arrhenius form and no measurement in
this work supplies the barriers, so the chain is switched off in the calculation
of \cref{sec:results} and its consequences are discussed in
\cref{sec:conclusions}.

\subsection{The size distribution and its closure}
\label{sec:moments}

Carrying only $n$ and $c$ assigns every loop at a point the same size, and the
population that produces is monodisperse by construction. The second content
$q^k_{sL}$ removes that restriction. Its Fokker--Planck part is
\begin{equation}
  \left.\dot q^k_{sL}\right|_{\mathrm{FP}}
  = 2\Pi^k_{sL}\Bigl[\;
      \underbrace{\Sigma^{\mathrm{net}}_k\,\bar m^k_{sL}\,\Delta_k^{3/8}}_{\text{drift}}
      \;+\;
      \underbrace{\tfrac12\,\Sigma^{\mathrm{gro}}_k\,\Delta_k^{-1/8}}_{\text{spreading}}
    \Bigr],
  \label{eq:q-FP}
\end{equation}
to which nucleation, the floor current and the basal transfers are appended
weighted by $m^{2}$ rather than $m^{1}$ or $m^{0}$:
\begin{equation}
  \dot q^k_{sL}=\left.\dot q^k_{sL}\right|_{\mathrm{FP}}
    + J^k_{sL}\bigl(m^{\mathrm{nuc}}_k\bigr)^{2}
    - m_{\min}^{2}\,\Phi^k_{\min}
    \;\pm\;\Gamma^{q}_{\theta}.
  \label{eq:qdot}
\end{equation}
The spreading term does not vanish when the family stops growing: at
$\Sigma^{\mathrm{net}}_k=0$ the drift term is zero and
$\dot q^k_{sL}=\Pi^k_{sL}\Sigma^{\mathrm{gro}}_k\Delta_k^{-1/8}>0$, so the
distribution keeps broadening at fixed mean size and fixed number. There is no
value of the two lower moments that expresses that, which is the argument for
carrying the third.

Closing \cref{eq:qdot} with a log-normal $f_k(m)$ determines its two parameters
from the three carried moments,
\begin{equation}
  \bar m^k_{sL}=\frac{c^k_{sL}}{n^k_{sL}},
  \qquad
  \Delta_k\equiv\frac{q^k_{sL}\,n^k_{sL}}{\bigl(c^k_{sL}\bigr)^{2}}=e^{s_k^{2}}\ \ge 1,
  \qquad
  s_k^{2}=\ln\Delta_k .
  \label{eq:lognormal}
\end{equation}
$\Delta_k\ge1$ is Cauchy--Schwarz on a non-negative measure, so
\textbf{$\Delta_k<1$ is not a loose tolerance but a defect}: it says the three
carried numbers are not the moments of any distribution. Keeping the triple
realizable places a real constraint on how each channel debits $q$, and
\cref{app:moments} states the three rules that follow --- of which the one most
easily got wrong is coalescence, because a coalescing loop does not leave the
family, it \emph{merges}, so fewer loops holding the same content are
\emph{larger} loops and coalescence \emph{raises} $q$. Debiting $q$ by the
removed number times $\langle m^{2}\rangle$, the natural-looking reading,
inverts the sign of the largest single term in the equation.

The closure reaches the rest of the model in exactly one place, as
$\Delta_k^{-1/8}$ on the sink prefactor, and it reaches the discrete conversion
of \cref{sec:conversion} as the distribution loop sizes are drawn from.

\subsection{Summary of the equations}
\label{sec:summary-equations}

The system integrated is two blocks with different structure, and the difference
between them is the whole of the numerical method.

\paragraph{Mobile block --- four fields, spatially coupled, no time derivative.}
\begin{equation}
  \spvec 0 = \nabla\!\cdot\!\bigl(\spmat D_{\mathrm{ten}}\nabla\spvec c_M\bigr)
    + \spvec P_M
    + \tfrac12(\spmat I\otimes\spvec c_M^{T})\spmat R_2\spvec c_M
    + (\spmat E-\spmat D)\spvec c_M + \spvec S_{vL},
  \qquad
  \spvec c_M\big|_{\partial\Omega_D}=\spvec c_M^{\mathrm{eq}},
  \label{eq:block-mobile}
\end{equation}
$4\times N_{\mathrm{node}}$ unknowns, nonlinear and algebraic, coupled across
the mesh by the divergence term alone.

\paragraph{Immobile block --- 27 fields, pointwise, no spatial coupling.}
For each family $(s,k)\in\mathcal F$ at each node independently,
\begin{align}
  \dot n^k_{sL} &= J^k_{sL} - \Delta N - \Phi^k_{\min} \pm \Gamma^{n}_{\theta}
                   - \nu_{gb}\,\Theta^{k}_{0},
  \label{eq:block-n}\\
  \dot c^k_{sL} &= \left.\dot c^k_{sL}\right|_{\mathrm{growth}}
                   + G_k - \Delta C \pm \Gamma^{c}_{\theta}
                   - \nu_{gb}\,\Theta^{k}_{1},
  \label{eq:block-c}\\
  \dot q^k_{sL} &= \left.\dot q^k_{sL}\right|_{\mathrm{FP}}
                   + J^k_{sL}(m^{\mathrm{nuc}}_k)^{2} - m^{2}_{\min}\Phi^k_{\min}
                   \pm \Gamma^{q}_{\theta} - \nu_{gb}\,\Theta^{k}_{2},
  \label{eq:block-q}
\end{align}
with $\Theta^k_j$ the size-gated grain-boundary removal of \cref{sec:gb-loops}
and $\Gamma_\theta$ the basal-chain transfers of \cref{eq:basal-chain}.
Together with the evolving network density $\rho_N$ and the atom-balance
accumulators of \cref{sec:results-conservation}, this is a system of ordinary
differential equations at each node --- $3\times9+1$ fields plus accumulators ---
in which no neighbor's state appears. \textbf{The mobile species are the only
thing that couples one node to another}, and freezing them, as
\cref{sec:split} does, removes the spatial coupling by construction rather
than by approximation.
